% =====================================================================
% Rozdział 5 – Testy i analiza wyników (~10 str.)
% Systematycznie: zbiór testowy, wyniki referencyjne, miary ilościowe,
% porównanie z wersją wyjściową, odniesienie do innych rozwiązań.
% =====================================================================
\chapter{Testy i~analiza wyników}
\label{ch:testy}

% ---------------------------------------------------------------------
\section{Metodyka testów}
\label{sec:metodyka}
% - Stanowisko testowe: JetBot z obudową, otoczenie.
% - Scenariusze: obiekt statyczny w różnych kierunkach i odległościach,
%   obiekt w obszarze nakładania, kilka obiektów, obrót robota, różne
%   oświetlenie.
% - Pomiary wydajności w firmware: średnie i maksima z bloków po
%   50 klatek, pierwszy blok odrzucany, jedna zmiana na krok, stałe
%   warunki.
% - Wzorzec prawdy dla kierunków: np. znaczniki kątów na podłodze.
% Stan: ✅ (wydajność), ⏳ (scenariusze systemowe)


% ---------------------------------------------------------------------
\section{Ocena modelu i~jakość detekcji}
\label{sec:jakosc-detekcji}
% - Metryki Edge Impulse na zbiorze testowym: F1, precision, recall,
%   macierz pomyłek.
% - Zasoby: RAM (arena ok. 270 kB), flash, szacowany vs zmierzony czas
%   inferencji na S3 (97 ms); int8 vs float32 (jeśli dostępne).
% - Pojedyncza kamera: skuteczność od odległości, położenia w kadrze
%   (środek vs krawędź), oświetlenia.
% - Cały system: wykrywanie w każdym kierunku 360°, także na granicach
%   kamer; fałszywe detekcje i pominięcia.
% - Porównanie kamer o mniejszym i większym FOV. ❓ jeśli zrobisz.
% - Tab.: metryki modelu.
% - Rys.: macierz pomyłek.
% - Rys.: skuteczność od odległości.
% Stan: ⏳


% ---------------------------------------------------------------------
\section{Dokładność kalibracji i~fuzji}
\label{sec:dokladnosc-kalibracji}
% - RMS reprojekcji wszystkich kamer (podsumowanie z
%   \ref{sec:kalibracja-wewn}).
% - Błąd azymutu: zmierzony vs rzeczywisty, dla wielu kierunków wokół
%   robota.
% - Metoda A vs B2 na tym samym zestawie – zwłaszcza przy krawędziach
%   kadru.
% - Skuteczność fuzji: obiekt w obszarze nakładania daje jeden wynik,
%   nie dwa.
% - Rys.: błąd azymutu w funkcji kierunku.
% - Tab.: A vs B2.
% Stan: ⏳


% ---------------------------------------------------------------------
\section{Wydajność i~opóźnienia}
\label{sec:wydajnosc}
% - Przepustowość (FPS) a opóźnienie (od zdjęcia do wyniku) – to nie to
%   samo.
% - Przebieg optymalizacji węzła (porównanie z wersją wyjściową):
%   1,8 -> 6,64 -> 7,08 -> 8,11 -> 9,13 -> 8,91 FPS;
%   opóźnienie 668 -> 171 ms.
%   Kroki: wyłączenie kontroli sterty, -O2 dla LUT, cache 64 kB, trwały
%   interpreter, XCLK 20 MHz. Krok 6 (wejście z uint8) – sprawdzony,
%   nie wszedł.
% - Budżet opóźnień węzła: 171 ms = wiek klatki 61 ms + LUT 4,5 ms
%   + DSP 5,3 ms + sieć neuronowa 97 ms + reszta; wiek klatki ≈ 1,5 okresu
%   kamery (GRAB_LATEST, 2 bufory). Rozkład: sieć ≈ 87%, LUT ≈ 4%.
% - Część po stronie komputera: transmisja Wi-Fi, odbiór, fuzja,
%   publikacja ROS 2. ⏳
% - Cały system: 6 węzłów naraz – FPS, opóźnienie end-to-end
%   (zdjęcie -> komunikat ROS 2), utrata pakietów UDP. ⏳
% - Znaczenie dla robota: przy obrocie 30°/s opóźnienie 171 ms ≈ 5° błędu
%   kierunku (rząd jednej komórki FOMO) – stąd znacznik czasu zdjęcia.
% - Tab.: przebieg optymalizacji.
% - Rys.: FPS i opóźnienie po krokach.
% - Rys.: oś czasu jednej klatki (Gantt).
% Stan: ✅ (węzeł), ⏳ (system)


% ---------------------------------------------------------------------
\section{Ograniczenia i~porównanie z~innymi rozwiązaniami}
\label{sec:ograniczenia}
% - Fałszywe detekcje (podobne obiekty, odbicia); wpływ oświetlenia.
% - Błąd kierunku przy ruchu robota – wpływ opóźnienia, skuteczność
%   kompensacji z odometrii.
% - Ograniczenia 96x96: małe/odległe obiekty, dokładność kątowa ≈ komórka
%   FOMO.
% - Ograniczenia Wi-Fi: zmienne opóźnienie, zakłócenia; pojedyncze błędy
%   sterownika kamery (EV-VSYNC-OVF), sendto() errno=12 przy starcie.
% - Porównanie z innymi rozwiązaniami (odniesienie do stanu wiedzy):
%   LiDAR, kamera 360°, kilka kamer na jednym komputerze; wyniki
%   FOMO/TinyML z literatury i benchmarków Espressif – koszt, FPS,
%   opóźnienie, dokładność kątowa.
% - Tab.: porównanie z innymi podejściami.
% Stan: ⏳

